% ===== §7 Developmental Dynamics under Relational Coupling =====
\section{Developmental Dynamics under Relational Coupling}
\label{p14:sec:dynamics}

\noindent
The previous section established that the relational field is a coupled dynamical system whose co-evolutionary response to contingent events is real, multi-level, and irreducible to the sum of individual dynamics. But the account so far has been primarily synchronic: it has described what happens at a given moment when a contingent event enters the relational field. The deeper question is diachronic: how does the relational system \emph{develop} over time? How do repeated contingent events---flowers and griefs, shared moments of beauty and shared moments of difficulty---accumulate into the trajectory of a relational life? What is the geometry of this trajectory, and how does it relate to the individual dynamics of the two people who compose the relational system?

This section addresses these diachronic questions in five movements, culminating in a formal account of event-driven relational dynamics drawn from the biology of spiking neural networks (SNN). The account is both a formal contribution---it provides a precise model of how contingent events trigger structural change in the relational system---and a philosophical one: it illuminates the temporal structure of relational flourishing in a way that connects the micro-level dynamics of shared moments to the macro-level trajectory of a shared life.

\subsection{The Problem of the Relational Phase Space}
\label{p14:subsec:phasespace}

\noindent
To speak of the \emph{development} of a relational system is to speak of its trajectory through a space of possible states. But what is the state space of a relational system? This question is more subtle than it appears, because the state space of the coupled system is not simply the Cartesian product of the state spaces of the two individuals.

Let $\mathcal{X}_1$ and $\mathcal{X}_2$ be the state spaces of the two individual dynamical systems---the spaces of all possible states of each individual, including their emotional states, attentional orientations, cognitive sets, and bodily dispositions. The state space of the coupled system is, at first approximation, the product space $\mathcal{X}_1 \times \mathcal{X}_2$. But this approximation misses something essential: the relational field has its own degrees of freedom---the history of shared events, the accumulated attunement, the quality of the \emph{between}---that are not reducible to the current states of the two individuals. The full state space of the relational system must include these relational degrees of freedom:

\begin{equation}
\mathcal{X}_{\mathrm{rel}} = \mathcal{X}_1 \times \mathcal{X}_2 \times \mathcal{R}
\end{equation}

where $\mathcal{R}$ is the space of relational states---the space of all possible configurations of the \emph{between}, including the history of shared events, the current quality of attunement, and the geometry of the relational field as shaped by past co-evolution. The dimension of $\mathcal{R}$ is, in general, much larger than the dimensions of $\mathcal{X}_1$ and $\mathcal{X}_2$: the relational field carries more information than either individual, because it carries the history of their interactions as well as their current states.

This extended state space is the proper arena for describing relational development. The trajectory of a relational system through $\mathcal{X}_{\mathrm{rel}}$ is the trajectory of the couple's shared life: each point on the trajectory corresponds to a moment in that life, with its particular configuration of individual states and relational states, and the trajectory as a whole is the geometric realisation of the shared life in the space of possible relational configurations.

\subsection{The Emergence and Stability of Shared Attractors}
\label{p14:subsec:attractors}

\noindent
Within this extended phase space, the co-evolutionary dynamics of the relational system produce attractors that are genuinely relational---attractors that exist in $\mathcal{X}_{\mathrm{rel}}$ but correspond to no attractor in either $\mathcal{X}_1$ or $\mathcal{X}_2$ alone. These shared attractors are the formal representation of what, in phenomenological terms, we call relational modes: the characteristic ways of being together that an intimate relation develops over time---the particular rhythm of shared attention, the characteristic register of shared affect, the habitual mode of joint response to contingent events.

A shared attractor is, formally, a region of $\mathcal{X}_{\mathrm{rel}}$ toward which the relational dynamics tend from a range of initial conditions. Its basin of attraction---the set of initial states from which the system converges to the attractor---is a measure of the attractor's stability: a large basin indicates a robust relational mode that is resilient to perturbation, while a small basin indicates a fragile mode that is easily displaced by contingent events or internal fluctuations.

The development of a relational system can be described, in these terms, as the progressive emergence and stabilisation of shared attractors. Early in a relation, the coupled dynamics are highly sensitive to initial conditions and to contingent events: small perturbations produce large divergences, and the system has not yet developed the stable relational modes that characterise a mature intimate relation. Over time, through repeated co-evolutionary responses to contingent events, shared attractors emerge and their basins of attraction expand: the relational system becomes more stable, more predictable in its characteristic modes, more resilient to perturbation.

This developmental trajectory is the formal correlate of what is experienced, from within the relation, as the deepening of intimacy: the sense that the relation has developed its own characteristic ways of being together, that two people have learned each other in a way that makes their joint responses to contingent events feel natural and effortless, that the \emph{between} has acquired a richness and depth that it did not have at the beginning. The emergence of shared attractors is the dynamical ground of this experienced deepening.

\begin{claim}[Shared attractors as relational achievements]
The development of a relational system consists, in part, in the emergence and stabilisation of shared attractors in the relational phase space $\mathcal{X}_{\mathrm{rel}}$. These attractors are genuine relational achievements: they do not exist in the state space of either individual and could not have been reached by either individual's dynamics alone. The deepening of intimacy, experienced phenomenologically as the growth of shared modes and mutual attunement, corresponds, at the dynamical level, to the expansion of the basins of attraction of these shared attractors.
\end{claim}

\subsection{The Nesting of Individual and Relational Dynamics}
\label{p14:subsec:nesting}

\noindent
A question that the dynamical account must address is the relationship between individual and relational dynamics: if the relational field has its own attractors, its own trajectory, its own degrees of freedom, what becomes of the individual? Is the individual dissolved into the relational field, their dynamics entirely subsumed by the coupled dynamics?

The answer, as anticipated in \xref{p14:sec:turn}, is no---but the relationship is more subtle than simple parallel coexistence. The individual dynamics and the relational dynamics are \emph{nested}: the individual dynamics are embedded within and shaped by the relational dynamics, but they retain a genuine autonomy within that embedding. This nesting can be made precise in the language of dynamical systems theory through the concept of \emph{slaving}: in a coupled system with multiple timescales, the fast variables are ``slaved'' to the slow variables, in the sense that they rapidly relax to a manifold determined by the slow variables \parencite{strogatz2001}.

In the relational system, the relational degrees of freedom $\mathcal{R}$---the accumulated history, the attractor landscape, the quality of attunement---change slowly relative to the rapid fluctuations of individual state. The individual dynamics $\mathcal{X}_1$ and $\mathcal{X}_2$ are, in this sense, slaved to the relational dynamics: they evolve rapidly within a landscape that is shaped and continuously reshaped by the slower evolution of the relational field. The individual dynamics are not eliminated by this nesting; they are the fast, fine-grained activity that occurs within the slow, coarse-grained structure of the relational field.

This nesting has several important consequences. First, it means that the individual retains genuine dynamical autonomy within the relational field: the fast fluctuations of individual state are not fully determined by the relational field, and the individual's responses to contingent events retain an element of genuine particularity---a specificity of response that is not simply the average of the relational field's response. Second, it means that the relational field is genuinely constraining: the individual dynamics do not unfold in a vacuum but within a landscape shaped by the relational history, and this landscape channels the individual dynamics in directions that are consistent with the relational attractor structure. Third, it means that the individual dynamics continuously feed back into the relational dynamics: the fast fluctuations of individual state are integrated, over time, into the slow evolution of the relational field, so that the individual's particular responses to contingent events gradually reshape the relational landscape that constrains future responses.

\subsection{The Time-Evolution of Coupling Strength}
\label{p14:subsec:coupling}

\noindent
The coupling strength $\varepsilon$ in the formal model of \xref{p14:subsubsec:coupled} is not a fixed parameter but a dynamic variable: it changes over time in response to the history of the relational system and to the contingent events that enter it. Understanding the time-evolution of coupling strength is essential for understanding the developmental trajectory of the relational system and the conditions under which relational happiness is possible.

Three mechanisms govern the time-evolution of coupling strength:

\emb{Baseline coupling through sustained co-presence.} The ordinary daily co-presence of an intimate couple---the sharing of space, time, attention, and activity---produces a sustained baseline level of coupling that is the ground condition of relational co-evolution. This baseline coupling is not dramatic; it is the quiet, continuous mutual influence of two people who live together, whose rhythms of sleep and waking, eating and working, are attuned to each other's. But it is real and dynamically significant: it maintains the relational system in a regime of moderate coupling from which it can respond to contingent events, and it gradually strengthens the shared attractors that constitute the relational landscape.

\emb{Event-driven coupling spikes.} Contingent events---flowers on the path, shared griefs, moments of surprise or delight---produce transient spikes in coupling strength. In the formal model, a contingent event $e$ occurring at time $t_e$ produces a transient increase in $\varepsilon$:

\begin{equation}
\varepsilon(t) = \varepsilon_0 + \sum_k \Delta\varepsilon_k \cdot h(t - t_k)
\end{equation}

where $\varepsilon_0$ is the baseline coupling, $\Delta\varepsilon_k$ is the magnitude of the coupling increase produced by event $k$, and $h(t)$ is a response kernel that captures the temporal profile of the coupling increase---typically a rapid rise followed by a slower decay. The happiness of the shared flower corresponds, in this model, to the transient state of enhanced coupling: the moment in which the relational system is operating in a regime of stronger-than-baseline coupling, producing more intense co-evolutionary dynamics and a richer shared attractor structure.

\emb{Long-term structural modification.} Some contingent events---particularly those that are intense, unexpected, or deeply significant---produce not merely transient coupling increases but permanent structural modifications of the relational landscape. These events reshape the attractor structure of the relational field in ways that persist long after the event itself has passed: they create new shared attractors, expand the basins of existing attractors, or, in cases of relational crisis, contract or destroy attractor basins. This permanent structural modification is the dynamical correlate of what is experienced as a turning point in a relation---a moment after which the relation is, in some fundamental sense, different from what it was before.

\subsection{Event-Driven Relational Dynamics: The SNN Analogy}
\label{p14:subsec:snn}

\noindent
The most powerful formal model for understanding event-driven structural modification in dynamical systems is the spiking neural network (SNN)---a class of computational models inspired by the biology of neural circuits, in which information is processed through discrete, temporally precise events (spikes) rather than through continuous rate codes \parencite{maass1997,gerstner1996}. The SNN model provides a precise and biologically grounded analogy for the way in which contingent events trigger structural change in the relational system.

\subsubsection{The SNN as a Biological Dynamical System}
\label{p14:subsubsec:snnbio}

\noindent
In a biological neural network, individual neurons integrate synaptic inputs over time and generate an action potential (spike) when their membrane potential exceeds a threshold. The spike propagates along the axon and triggers the release of neurotransmitters at synaptic terminals, which in turn influence the membrane potentials of postsynaptic neurons. The network's behaviour is thus determined by the pattern of spikes---their timing, their spatial distribution across the network, and the synaptic weights that determine the influence of each spike on its targets.

The crucial property of SNNs that motivates the relational analogy is \emph{spike-timing dependent plasticity} (STDP): the synaptic weights of the network are modified by experience, and the direction and magnitude of this modification depend on the precise temporal relationship between pre- and postsynaptic spikes \parencite{bi1998}. Specifically:

\begin{equation}
\Delta w_{ij} = \begin{cases}
A_+ \exp\!\left(-\dfrac{t_{\mathrm{post}} - t_{\mathrm{pre}}}{\tau_+}\right) & \text{if } t_{\mathrm{post}} > t_{\mathrm{pre}} \\[6pt]
-A_- \exp\!\left(-\dfrac{t_{\mathrm{pre}} - t_{\mathrm{post}}}{\tau_-}\right) & \text{if } t_{\mathrm{post}} < t_{\mathrm{pre}}
\end{cases}
\end{equation}

where $\Delta w_{ij}$ is the change in synaptic weight from neuron $j$ to neuron $i$, $t_{\mathrm{pre}}$ and $t_{\mathrm{post}}$ are the times of the pre- and postsynaptic spikes, $A_+$ and $A_-$ are the magnitudes of potentiation and depression, and $\tau_+$ and $\tau_-$ are the time constants of the learning windows. The key biological insight is that the causal order of spikes matters: if the presynaptic neuron fires before the postsynaptic neuron (which suggests that the presynaptic neuron contributed to causing the postsynaptic spike), the synapse is strengthened; if the order is reversed, the synapse is weakened.

STDP is Hebbian learning at millisecond precision: it implements, at the neural level, the principle that ``neurons that fire together, wire together''---but with the crucial temporal specificity that distinguishes genuine causal coupling from mere coincidence.

\subsubsection{The Relational Analogy}
\label{p14:subsubsec:analogy}

\noindent
The SNN model maps onto the relational system with a precision that is, we believe, philosophically illuminating. The mapping proceeds as follows:

\emb{The two individuals as coupled neural populations.} Each person in the relational system corresponds to a population of neurons (or, at the systems level, to a dynamical system with SNN-like properties): a complex, high-dimensional system that integrates inputs over time, generates event-like responses (emotional reactions, attentional shifts, affective spikes), and transmits these responses to the other system through the coupling of the relational field.

\emb{The relational field as the synaptic network.} The coupling between the two individual systems is mediated by the relational field---the \emph{between} with its accumulated history, its attunement, its shared attractor structure. This coupling corresponds, in the SNN analogy, to the synaptic network that connects the two neural populations: a structured set of connection weights that determines how the activity of one population influences the other, and that is itself modified by experience through the STDP-like mechanism of relational co-evolution.

\emb{Contingent events as relational spikes.} A contingent event---the flower on the path, the unexpected news, the moment of shared laughter---corresponds to a spike: a discrete, temporally localised event that perturbs the state of the relational system and propagates through the coupling to affect both individuals. The spike is not a gradual influence but a sudden perturbation that crosses a threshold and initiates a rapid, discrete response.

The threshold character of relational spikes is philosophically significant. Not every contingent occurrence that takes place in the vicinity of the relational field enters that field as a relational event; only those that are sufficiently salient, sufficiently surprising, or sufficiently resonant with the current state of the relational field cross the threshold of joint attention and become relational spikes. This threshold structure explains why the same environmental event---a flower on the same path---can enter the relational field as a vivid relational event on one occasion and pass unnoticed on another: the threshold depends on the current state of the relational system, including the quality of co-present attention and the degree of mutual attunement.

\emb{STDP as relational structural modification.} The most important element of the analogy is the mapping between STDP and the long-term structural modification of the relational field by contingent events. In the SNN, STDP modifies synaptic weights in a way that depends on the precise temporal relationship between spikes: spikes that occur in causal sequence strengthen the connections between the neurons that generated them, while spikes that occur in anti-causal sequence weaken those connections. In the relational system, a contingent event that is shared in a timely fashion---that enters the relational field through an act of sharing that occurs close in time to the event itself---produces a stronger structural modification of the relational field than an event that is shared with delay.

This temporal dependence of relational STDP provides a formal account of the phenomenological observation noted in \xref{p14:subsec:sharing}: the happiness of the shared flower depends crucially on the timeliness of the sharing. The flower that is shared immediately---``look!''---has a stronger co-evolutionary effect on the relational field than the flower that is mentioned later that evening: ``oh, I saw the most beautiful flower today.'' The later mention is still a form of sharing, and it still enters the relational field as a relational event; but the temporal gap between the event and the sharing reduces the magnitude of the relational STDP effect, producing a weaker structural modification of the relational landscape.

Formally, let $t_e$ be the time of the contingent event and $t_s$ be the time of the sharing act. The magnitude of the structural modification is:

\begin{equation}
\Delta \mathcal{R} \propto A \cdot \exp\!\left(-\frac{t_s - t_e}{\tau_{\mathrm{rel}}}\right)
\end{equation}

where $\tau_{\mathrm{rel}}$ is the time constant of relational STDP---the characteristic timescale over which the sharing must occur for the contingent event to produce its maximal co-evolutionary effect. Events shared within this window produce strong structural modification; events shared outside it produce progressively weaker modification; events that are never shared produce no structural modification of the relational field, however significant they may be for the individual who experienced them.

\begin{claim}[Relational STDP and the timeliness of sharing]
The structural modification of the relational field produced by a contingent event depends on the temporal proximity of the sharing act to the event itself. Events shared immediately produce stronger co-evolutionary effects---stronger modifications of the relational attractor landscape---than events shared with delay. This relational STDP principle provides a formal account of the phenomenological observation that the happiness of the shared flower depends on the timeliness of the sharing, and it implies a practical principle: the cultivation of relational happiness requires the development of habits of timely sharing, of the kind of co-present attentiveness that makes immediate sharing possible.
\end{claim}

\subsubsection{The Threshold Mechanism and Relational Readiness}
\label{p14:subsubsec:threshold2}

\noindent
The SNN analogy also illuminates the threshold mechanism introduced in \xref{p14:subsec:threshold}. In the SNN, a neuron generates a spike only when its membrane potential exceeds a threshold: subthreshold inputs accumulate but do not produce a discrete response, while suprathreshold inputs trigger the rapid, all-or-nothing response of the action potential. This threshold mechanism is essential to the SNN's information-processing properties: it ensures that only sufficiently strong or sufficiently coincident inputs produce a discrete response, filtering out weak or uncorrelated signals.

In the relational system, the threshold for the generation of a relational spike---for the entry of a contingent event into the relational field as a relational event---is determined by the current state of the relational field. A field in a state of high co-present attunement has a lower threshold: smaller and more subtle contingent events can enter the field as relational spikes. A field in a state of low attunement, distraction, or emotional disconnection has a higher threshold: only large, dramatic, or emotionally intense events can penetrate the threshold and enter the field as relational spikes.

This threshold modulation by relational state explains a familiar phenomenological observation: the same environment---the same path, the same flowers, the same afternoon light---feels more vivid, more alive, more generative of shared happiness when one is in a state of genuine co-present attunement with a beloved person than when one is distracted, preoccupied, or emotionally absent. The environment has not changed; the threshold has changed. The practices of relational attentiveness that \xref{p14:sec:praxis} will discuss are, in part, practices of threshold modulation: they cultivate the state of co-present attunement that lowers the threshold for relational spikes and thereby increases the relational field's sensitivity to contingent events.

\subsubsection{From Spikes to Trajectories: The Cumulative Structure of Relational Development}
\label{p14:subsubsec:cumulative}

\noindent
The SNN analogy, applied to the relational system, yields a picture of relational development as the cumulative effect of a lifetime of relational spikes. Each contingent event that enters the relational field as a relational spike produces a small structural modification of the relational landscape---a tiny shift in the attractor structure, a small change in the coupling weights, a modest expansion or contraction of a basin of attraction. Over the course of a shared life, these small modifications accumulate into the rich, complex relational landscape of a mature intimate relation: its characteristic modes of joint attention, its habitual registers of shared affect, its resilience in the face of difficulty, its capacity for the deepest forms of co-present happiness.

This cumulative structure is the dynamical correlate of what \pinkref{Paper~IX} described in terms of geometric phase and holonomy \parencite{paperix}: the spiral structure of relational development, in which each traversal of the relational cycle returns not to the starting point but to a new point carrying the accumulated phase of the journey. The relational STDP mechanism provides the micro-level account of how this accumulation occurs: each shared contingent event produces a small modification of the relational landscape, and the integral of these modifications over time is the holonomy of the relational trajectory---the total accumulated phase that constitutes the spiral of a shared life.

The happiness that arises in the mature intimate relation---the deep, quiet happiness of two people who have shared a long history of contingent events, who have developed the shared attractor landscape of a life together, whose coupling is strong and whose threshold is low---is not the same as the happiness of the early relation, with its dramatic spikes and its rapidly changing landscape. It is a happiness of depth rather than of intensity: the happiness of a relational system that has accumulated, through the patient work of a lifetime of sharing, a rich and stable landscape of shared attractors, a deep and resilient coupling, a \emph{between} that is dense with the history of everything that has been received together.

\begin{claim}[The cumulative structure of relational happiness]
The happiness of a mature intimate relation is not the sum of individual moments of shared happiness but the cumulative product of a lifetime of relational STDP: the rich attractor landscape that emerges from the structural modifications produced by a lifetime of shared contingent events. This cumulative happiness has a different character from the happiness of individual shared moments---it is quieter, deeper, more stable, and more resilient---and it constitutes the fullest realisation of what the GRB framework calls relational \emph{eudaimonia}: the flourishing of a relational system that has developed, through sustained co-evolution, the depth and stability of a genuinely shared life.
\end{claim}

With this formal account of relational developmental dynamics in place, we are ready to examine how these dynamics can be studied empirically. The next section proposes a research methodology for testing the claims of the GRB account---a methodology that must be as attentive to the distinctive ethical challenges of relational research as it is to the formal requirements of empirical science.
